\subsection{Matriz de evaluación ponderada}

Con base en los criterios anteriores, se aplicó una matriz de decisión ponderada en una
escala discreta de 1 a 5, donde 5 representa el mayor grado de adecuación al objetivo
del trabajo y 1 el menor. Los puntajes corresponden a una valoración relativa respecto a
la posibilidad de integrar y evaluar una estrategia de compresión sin pérdida sobre el
vector de estado.

La asignación de valores se estableció a partir de dos fuentes complementarias de
evidencia: la documentación oficial y literatura técnica asociada a cada proyecto, y una
revisión directa del código disponible en los repositorios o artefactos accesibles de
cada simulador. La documentación permitió identificar el propósito de diseño, el
modelo de ejecución y las capacidades de paralelismo reportadas por cada herramienta,
mientras que la inspección del código permitió valorar la localización de la estructura
que representa el estado, el nivel de acoplamiento entre esa estructura y la lógica de
simulación, y la facilidad de introducir modificaciones sobre el backend de memoria.
Bajo esta escala, la evaluación obtenida fue la siguiente:

\begin{table}[H]
\centering
\caption{Evaluación de alternativas (escala 1--5) y total ponderado}
\label{tab:evaluacion_sim}
\begin{tabular}{lcccc}
\hline
\textbf{Simulador} & \textbf{C1} & \textbf{C2} & \textbf{C3} & \textbf{Total ponderado} \\
\hline
QuEST          & 2 & 2 & 5 & 2.60 \\
Intel-QS       & 2 & 2 & 5 & 2.60 \\
qsim           & 4 & 4 & 4 & 4.00 \\
Quantum++      & 4 & 4 & 2 & 3.60 \\
TMFQSfullstate & 5 & 5 & 3 & 4.60 \\
\hline
\end{tabular}
\end{table}

Los valores consignados en la tabla responden a la forma en que cada alternativa
articula flexibilidad arquitectónica y contexto de ejecución. QuEST e Intel-QS
obtuvieron la máxima valoración en proyección de escalabilidad debido a su orientación
explícita hacia ejecución multinúcleo, distribuida y, en el caso de QuEST, acelerada
por GPU \cite{Jones2019, quest_repo, intelqs_doc}. Sin embargo, tanto la
documentación como la revisión del código muestran una arquitectura más orientada a
desempeño sobre infraestructuras HPC que a modificación experimental del backend de
memoria, lo que reduce su valoración en los dos primeros criterios.

En qsim, la separación entre componentes de simulación y manejo del \emph{state
vector}, visible tanto en su documentación como en la organización del código, ofrece
una base más modular para intervenir el almacenamiento con cambios controlados
\cite{qsim_repo, qsim_cirq}. Por ello, recibió valores altos en intervenibilidad y
desacoplamiento, además de una valoración también alta en escalabilidad por sus
variantes orientadas a ejecución eficiente.

Quantum++ obtuvo una valoración alta en los dos primeros criterios por su diseño como
biblioteca general en C++ de estructura ligera, distribuida como librería
\emph{header-only}, lo que favorece la inspección, modificación y extensión del código
\cite{qpp_paper}. No obstante, su proyección de escalabilidad resulta menor que la de
herramientas concebidas desde el inicio para contextos HPC o ejecución distribuida,
razón por la cual su valoración en el tercer criterio fue inferior.

TMFQSfullstate alcanzó la valoración más alta en intervenibilidad del almacenamiento y
desacoplamiento entre lógica de simulación y backend de memoria. Tanto la literatura,
como la revisión directa del código muestran una base de implementación más
directa, con una organización favorable para sustituir la representación del estado y
mantener control sobre la lógica de operación del simulador
\cite{diaz2024software, diaz_mdpi_mmss, tmfqs_repo}. Aunque su contexto de
escalabilidad es más limitado que el de QuEST o Intel-QS, conserva una proyección
suficiente para desarrollos posteriores y, sobre todo, ofrece la mejor adecuación
metodológica para integrar la estrategia de compresión propuesta.
